\section{Discussion and Limitations}
\label{sec:discussion}

\subsection{Interpretation of Results}

Taken together, Sections~\ref{sec:formal} and~\ref{sec:evaluation} support
the paper's central claim, but with an asymmetry worth stating plainly: the
empirical evaluation runs the full, composed, real system (all layers
integrated, real Go code, real infrastructure), while the completed formal
verification is currently per-layer only -- each layer shown internally
free of violations in isolation, not yet shown to correctly refine the
abstract specification as a composed whole (Section~\ref{sec:formal}). This
means the live testbed currently carries more of the burden of demonstrating
whole-system correctness than the formal verification does; the formal side
contributes a precise statement of the target safety-and-accountability
property, machine-checked evidence at the component level, and a
five-mechanism ablation study showing each mechanism is individually
necessary at the FSM layer, while the live evaluation shows the resulting
design behaves correctly under real, not merely modeled, faults -- most
directly in the RQ2 composite-fault trace, where a real
\textsc{Sovereign}$\to$\textsc{Recovering}$\to$\textsc{Anchored} cycle
completes with 100\% AppHash agreement under three overlapping fault classes
at once. Neither result alone would be sufficient: the TLA+ model does not
capture real network timing, container resource contention, or the actual
Noir/Barretenberg proving pipeline, while the live testbed alone cannot
explore the state space exhaustively enough to guarantee the absence of
subtle safety violations, and its qualitative demonstrations (composite
fault, seven-scenario sweep, the P4 cascade, the reorg-depth reaction) are
each a single trial rather than a repeated, statistically characterized
sample -- unlike the properly repeated E5 hysteresis sweep, the E7 overhead
measurement ($n{=}5$), and the E4 synthetic Monte Carlo study (2,000
trials/cell). We report exact transition timings from these single-trial
runs (e.g.\ ``$t{=}321$\,s'') as one observed instance of correct behavior,
not as a characterized distribution. Repeating these live trials is
straightforward future work; the cross-layer refinement check discussed in
Section~\ref{sec:formal} is a separate, open-ended gap whose resolution we
cannot schedule (Section~\ref{sec:conclusion}).

\textbf{Open question: asymmetric hardening.} Our RQ3 results
(Section~\ref{sec:evaluation}) surface a genuine, unresolved design gap.
Exponential backoff exists on the \textsc{Recovering}$\to$\textsc{Sovereign}
edge specifically to defend against a cheap, repeatable attack that forces
premature demotion, but the same attack shape is available on the
\textsc{Anchored}$\leftrightarrow$\textsc{Suspicious} boundary, whose
thresholds do not grow the more they are triggered. A partial mitigation
exists at the application layer (a queue cap on forced-transaction admission
while \textsc{Suspicious}), but whether this edge needs the same backoff
hardening already applied to \textsc{Sovereign}/\textsc{Recovering} remains
open, and we flag it as a concrete direction for follow-up work
(Section~\ref{sec:conclusion}) rather than resolving it speculatively here.

\textbf{Griefing surface created by the fallback itself.} Forcing
\textsc{Sovereign} is not only a liveness question but a security
tradeoff: an adversary who degrades only a validator-visible \emph{view}
of Bitcoin, DA, or P2P health -- without compromising any validator's
keys -- can drive a supermajority-visible cascade
(Section~\ref{sec:evaluation}, P4) and lock withdrawals network-wide for
as long as the degraded view persists. We have not quantified the
attacker's cost against this benefit (e.g.\ the cost of sustaining a
network-level attack against a supermajority of validators' Bitcoin/DA/P2P
connectivity, versus the value of a network-wide withdrawal freeze to that
attacker), nor analyzed the security implication of a long-range fork
between the Bitcoin-anchored history and the sovereign history at
re-anchoring time. Both are open problems this paper's evaluation does not
resolve.

\subsection{Limitations}
\label{subsec:limitations}

Section~\ref{subsec:sim-to-real} already discloses three gaps between this
testbed and a genuinely distributed, mainnet-scale deployment: a shared
physical host, an accelerated Bitcoin timescale, and single-source-per-validator
DA trust. Rather than restate those facts, we state here what each one
means for how the results in Section~\ref{sec:evaluation} should be read,
plus three further limitations not covered there.

\textbf{Single physical host.} Because every validator, plus the Bitcoin
and Celestia networks, share one machine's resources, our measured
overhead numbers (RQ4) are a reasonable but not definitive proxy for a
genuinely distributed deployment, and our wide-area-network fault
injection, while using kernel-level traffic control, approximates rather
than reproduces true geographic distribution.

\textbf{Accelerated Bitcoin timescale.} Because the testbed's Bitcoin
network runs roughly $30\times$ faster than mainnet, readers should treat
our wall-clock timings (e.g.\ the 341-second RQ2 trace) as demonstrating
\emph{correct sequencing} of the recovery process, not as mainnet-calibrated
latency figures; height-based thresholds are unaffected.

\textbf{Single-source-per-validator DA trust.} Because each validator
trusts only its own bridge rather than sampling across the other six, our
DA-outage results should be read as measuring a bridge (or all seven
together) going down, not as a sampling-detected partial failure of an
otherwise-healthy network; closing this gap is future work
(Section~\ref{sec:conclusion}).

\textbf{Single Byzantine actor at a time.} Every live attack in this paper
-- the malicious proposer, the ten-scenario suite in
Appendix~\ref{app:attacks}, and P3/P4 -- compromises exactly one validator
or attacker container at a time. The formal threat model bounds up to $f$
simultaneously Byzantine validators ($f{=}2$ at $N{=}7$), but no live
experiment has exercised $f$ validators colluding at once (e.g.\ two
validators jointly forging a certificate or double-signing in coordination);
single-actor tests establish a safety margin, not the formally-bounded
worst case. Relatedly, the ten-attack suite tests a Byzantine
\emph{proposer} forging its proposal's FSM-state field, but does not
separately test a Byzantine validator submitting a false individual health
certificate into the health-certificate aggregation step that precedes
consensus -- a related but distinct vector we have not confirmed is
covered by an existing test.

\textbf{Near-quorum split views untested.} Section~\ref{subsec:threat-model}
notes that validators running identical thresholds can still disagree if
their own external reads of Bitcoin, DA, or P2P health genuinely differ.
Every split view we tested is decisive (a self-healing 1-of-7 view in S5,
a sub-quorum 2-of-7 view in P3), not a persistent split straddling the
$\geq 2f{+}1$ quorum boundary (e.g.\ 3-of-7 or 4-of-7), which could in
principle cause repeated nil-prevotes and round failures rather than a
clean commit either way. We note this because it offers one plausible,
unconfirmed explanation for the $\sim$23--27\% block-production-rate
reduction observed even on the fully healthy RQ4 path
(Table~\ref{tab:overhead}): we describe that gap as ``round-skip
overhead'' without having isolated whether occasional round-level
disagreement of this kind, rather than pure validation latency, is what
drives it.

\textbf{Throughput tested only up to 100 tx/s.} RQ4's ``no throughput
ceiling'' finding (Section~\ref{sec:evaluation}) holds within the tested
range; we did not search for a rate at which the extended proposal's
per-block byte overhead compounds with mempool or gossip costs into a
measurable ceiling, so readers should not read this as a claim that no such
ceiling exists at higher throughput.

Beyond these six, the two threat-model boundaries introduced in
Section~\ref{subsec:threat-model} (the Bitcoin reorg-depth bound and the
peripheral-view convergence assumption) remain conditions on, rather than
caveats against, the safety and accountability property targeted in
Section~\ref{sec:formal}; we do not restate them a third time here.

\subsection{Threats to Validity}

\textbf{Internal validity.} Sharing one physical host (above) is the primary
threat to internal validity for our cost measurements; we mitigate this by
reporting both a ramped-load sweep and a fixed-rate, $n{=}5$-repeat
measurement with confidence intervals (Section~\ref{subsec:rq4}),
which shows the measured overhead gap is stable and reproducible rather than
noise, but does not eliminate host-sharing as a contributing factor to the
absolute numbers.

\textbf{External validity.} Results at 7-validator scale, on a single
physical host, may not extrapolate linearly to mainnet-scale validator sets
or genuinely distributed infrastructure; we regard our results as
establishing correctness and a bounded-overhead \emph{trend}, not as
mainnet-calibrated performance figures.

\textbf{Construct validity.} We use AppHash agreement and divergence
counts as our operational proxy for safety, and block
height progression / commit rate as our operational proxy for liveness,
consistent with how CometBFT-family systems are conventionally evaluated;
readers should note that this proxies the properties formally verified in
Section~\ref{sec:formal} rather than independently re-deriving them.